\begin{abstract}
עבודה זו בוחנת את השימוש ב־\tech{Gaussian Hidden Markov Model (HMM)} בשתי רמות: חיזוי כיוון התשואה ביום הבא, וזיהוי משטרי שוק סמויים בעלי משמעות סטטיסטית. הניסוי הסופי כולל תשעה נכסים מייצגים — \asset{SPY}, \asset{QQQ}, \asset{IWM}, \asset{TLT}, \asset{GLD}, \asset{HYG}, \asset{BTC-USD}, \asset{JPM} ו־\asset{NVDA} — בטווח \datecode{2014-01-01--2024-12-31}. לכל יום מחושבים ארבעה Features: תשואה לוגריתמית, תנודתיות מתגלגלת ל־20 ימים, טווח תוך־יומי ושינוי לוגריתמי בנפח. החלוקה כרונולוגית ומוגנת מזליגה: Train, Validation ו־Test, כאשר preprocessing מותאם על Train בלבד.

לכל נכס נבדקו \(K\in\{2,3,4\}\) וחמישה seeds. בחירת \(K\) ו־seed נעשתה לפי Validation log-likelihood בלבד. בנוסף ל־hard state assignment נבדקה תחזית המבוססת על כל וקטור ה־posterior. המודל הושווה ל־Baselines פשוטים ולשלושה מודלי Supervised ML: Logistic Regression, Random Forest ו־HistGradientBoosting.

מבחינת חיזוי, לא נמצא יתרון עקבי ל־HMM: ממוצע ה־DPA בין תשעת הנכסים היה \pct{54.54} ל־Naive train-mean, \pct{54.14} ל־Logistic Regression, \pct{53.27} ל־Random Forest, \pct{53.06} ל־HMM soft-posterior, \pct{52.49} ל־HistGradientBoosting ו־\pct{52.06} ל־HMM hard-state. לעומת זאת, ברמת ייצוג מבנה השוק התקבלו ממצאים חזקים יותר. ב־SPY, מודל \(K=4\) יצר states בעלי תשואה, volatility, drawdown, range ומשך שהייה שונים מאוד; ה־posterior entropy עלה בצורה חדה סביב מעברי state; מבנה המצבים היה יציב מאוד בין seeds; ומדד VIX, שלא שימש באימון, הפריד בין חלק מהמשטרים. בנוסף נמצא כי הקורלציות בין SPY לנכסים אחרים משתנות כתלות במשטר.

המסקנה המרכזית היא כי בפרויקט זה ה־HMM שימושי יותר כ־\textbf{latent market-state estimator} מאשר כמודל point forecasting ליום הבא. אין בעבודה טענה לרווחיות מסחר או ליתרון סטטיסטי כללי.
\end{abstract}

\textbf{מילות מפתח:} \tech{Hidden Markov Model}, משטרי שוק, \tech{posterior uncertainty}, \tech{directional prediction}, \tech{walk-forward}, \tech{Supervised Learning}, \tech{Reinforcement Learning}.
